\section{Why caching a mid-layer hidden works}
\label{sec:mechanism}
\paragraph{Division of labor.} Probing shows semantic-task accuracy peaks in
mid layers and falls at the top, while next-token ability forms only in the top
layers (consistent across Qwen and Llama). Causally, using only the front
$\sim$4--8 layers (depth 0.12--0.22) reaches 95\% of the full model's
downstream semantics. The cacheable depth ceiling ($j{\le}9$ zero-shot, pushed
to $j{\ge}12$ by self-distillation) coincides with this semantic-saturation
point.
\paragraph{Caching the top fails.} A negative control (caching a top-layer
hidden) collapses: the top hidden is a query-sensitive, pre-readout
representation, so caching it discards query conditioning
(qa5 $(12,0){=}61/50 \to (12,6){=}29/20$). Only caching the query-insensitive
bottom-$j$ hidden is safe---validating the front-$j$ design.
\paragraph{Causal end-to-end test.} Pretraining a model with a low-rank funnel
at layer $j$ (which forces the depth-$j$ hidden to be compressible) yields a
representation that degrades \emph{less} under rank-$r$ PCA compression of the
cache than a vanilla model at every rank, and the advantage \emph{grows with
model size} (1B$\to$3B $\approx 2\times$). This is direct evidence that a
compressible mid-depth hidden---not an incidental property---is what the
readout exploits.
